% Λύση του 29ου προβλήματος της ενότητας «Άλυτα Προβλήματα».
\subsection*{Λύση Προβλήματος 29 --- \en{SCMAGLEV} και αεροδυναμική ισχύς}

\begin{alist}
  \item Επειδή $P(v)=D(v)v$, για τις δύο σχεδιάσεις έχουμε
        \[
          {P_0(v)=kv^3}
          \qquad\text{και}\qquad
          {P_1(v)=0{,}87kv^3}.
        \]
        Για $v>0$ ισχύει
        \[
          P_0'(v)=3kv^2>0,
          \qquad
          P_1'(v)=2{,}61kv^2>0,
        \]
        άρα και οι δύο συναρτήσεις είναι γνησίως αύξουσες. Επίσης,
        \[
          P_0''(v)=6kv>0,
          \qquad
          P_1''(v)=5{,}22kv>0.
        \]
        Επομένως, είναι και οι δύο κυρτές στο $(0,+\infty)$.

        Η κυρτότητα σημαίνει ότι η απαιτούμενη πρόσθετη ισχύς για την
        ίδια αύξηση της ταχύτητας γίνεται όλο και μεγαλύτερη. Αυτό
        συμφωνεί με την κυβική εξάρτηση $P\propto v^3$.

  \item Επειδή χρησιμοποιούμε την ίδια μονάδα ταχύτητας στον αριθμητή
        και στον παρονομαστή,
        \begin{align*}
          \frac{P_1(603)}{P_0(500)}
          &=
          \frac{0{,}87k\cdot603^3}{k\cdot500^3}\\
          &=0{,}87\left(\frac{603}{500}\right)^3\\
          &\approx1{,}526.
        \end{align*}
        Άρα
        \[
          {P_1(603)\approx1{,}526P_0(500)}.
        \]
        Παρά τη μείωση της αντίστασης κατά $13\%$, στα
        $603\si{km/h}$ απαιτείται περίπου $52{,}6\%$ περισσότερη
        αεροδυναμική ισχύς από όση απαιτούσε η αρχική σχεδίαση στα
        $500\si{km/h}$. Η αιτία είναι ότι η ισχύς αυξάνεται με τον
        κύβο της ταχύτητας.

  \item Ζητούμε $P_1(v_*)=P_0(500)$. Επομένως,
        \begin{align*}
          0{,}87kv_*^3&=k\cdot500^3,\\
          v_*^3&=\frac{500^3}{0{,}87},\\
          v_*&=500\sqrt[3]{\frac{1}{0{,}87}}.
        \end{align*}
        Άρα
        \[
          {
            v_*=500\sqrt[3]{\frac{100}{87}}
            \approx523{,}8\si{km/h}
          }.
        \]
        Δηλαδή η αεροδυναμική βελτίωση επιτρέπει αύξηση της ταχύτητας
        κατά περίπου $23{,}8\si{km/h}$, ή $4{,}8\%$, χωρίς αύξηση της
        αεροδυναμικής ισχύος σε σχέση με την αρχική σχεδίαση στα
        $500\si{km/h}$.
\end{alist}
